% !TeX root = Asparsi.tex
% !BIB TS-program = biber
% !TeX encoding = UTF-8
% !TeX spellcheck = it_IT
\section{Esagono}
\begin{figure}
	\centering
	\includestandalone{geometria/esagono}
	\caption{Esagono}
	\label{fig:esagono1}
\end{figure}
\begin{thm}[Diagonali]
	Un esagono ha nove diagonali.
\end{thm}\index{Esagono!diagonale}\index{Diagonale!esagono}

\begin{figure}
	\centering
	\includestandalone{geometria/esagono_diagonali}
	\caption{Esagono diagonali}
	\label{fig:esagonodiagonali}
\end{figure}
\begin{proof}
	Consideriamo la~\vref{fig:esagonodiagonali}. Da ogni vertice escono tre diagonali i vertici sono sei quindi abbiamo \num{18} diagonali ma ogni diagonale collega due vertici quindi il numero totale delle diagonali è nove. 
	
	Come volevasi dimostrare.
\end{proof}
\begin{thm}[Lato esagono]\label{thm:latoesagonoinscritto}
	L'esagono inscritto in una circonferenza ha i lati uguali al raggio della circonferenza.
\end{thm}\index{Esagono!lato}
\begin{proof}
	Consideriamo la~\vref{fig:esagono1}. Il triangolo $OCD$ è isoscele di lati uguali $OC$ e $OD$. Quindi gli angoli in $C$ e in $D$ sono uguali. L'angolo in $O$ è la sesta parte di un angolo giro. Quindi
	\begin{align*} 
		2x+\ang{60}=&\ang{180}\\
		x=&\ang{60}
	\end{align*} 
	Dato che un triangolo equiangolo è equilatero i tre lati sono uguali.
	Come volevasi dimostrare.
\end{proof}
\begin{cor}[Lato triangolo equilatero]\label{cor:latotriangoloequilatero1}
	Dato un triangolo equilatero di lato $l$ inscritto in una circonferenza si raggio $r$ avremo  \[l=r\sqrt{3}\]
\end{cor}\index{Triangolo!equilatero}\index{Triangolo!equilatero!lato}
\begin{figure}
	\centering
	\includestandalone{geometria/esagono_triangolo_equilatero}
	\caption{Lato triangolo equilatero inscritto}
	\label{fig:latotriangoloequilatero1}
\end{figure}
\begin{proof}
	In una circonferenza di centro $O$ e raggio $r$ si disegna un esagono di lato $m$. Unendo fra loro i vertici non consecutivi otteniamo un triangolo come nella~\vref{fig:latotriangoloequilatero1}. 
	
	Infatti Consideriamo i triangoli $ABC$, $CDE$ e $EGA$. Questi poligoni hanno, per costruzione due lati e l'angolo fra essi compreso uguali e quindi i tre triangoli sono congruenti. Ne segue che il triangolo $CAE$ è equilatero.   
	
	Tracciamo il triangolo $EDA$. Dato che il vertice $E$ è sulla circonferenza e $DA$ è il diametro il triangolo è rettangolo. Per il~\vref{thm:latoesagonoinscritto} il cateto $DE$ ha lungnezza uguale al raggio.
	\begin{align*}
		EA^2=&DA^2-DE^2\\
		=&(2r)^2-r^2\\
		\intertext{Quindi}
		EA=&\sqrt{3r^2}\\
		=&\sqrt{3}r
	\end{align*}
	Da cui la tesi.
\end{proof}
\begin{thm}[Area esagono]\label{thm:areaesagonoin1}
	L'area dell'esagono di lato $R$ è \[A=\dfrac{3}{2}R^2\sqrt{3}\]
\end{thm}\index{Esagono!area}
\begin{proof}
	Consideriamo la~\vref{fig:esagono1}. Calcoliamo l'area del triangolo $OCD$. 
	\begin{align*}
		A_{OCD}=&Ra\dfrac{1}{2}
		\intertext{Abbiamo:}
		a=&\sqrt{\dfrac{4R^2-R^2}{4}}\\
		=&	\sqrt{\dfrac{4R^2-R^2}{4}}\\
		=&\dfrac{R}{2}\sqrt{3}\\
		\intertext{sostituendo}
		A_{OCD}=&\dfrac{R^2}{4}\sqrt{3}\\
		\intertext{L'esagono è formato da sei triangoli quindi:}
		A=&6\dfrac{R^2}{4}\sqrt{3}\\
		=&\dfrac{3}{2}R^2\sqrt{3}
	\end{align*}
	Come volevasi dimostrare.
\end{proof}
\begin{thm}[Rapporto tra aree]
	L'area di un triangolo equilatero inscritto in un esagono è la metà dell'area dell'esagono.
\end{thm}\index{Triangolo!equilatero!area}\index{Esagono!area}
\begin{proof}
	Consideriamo la~\vref{fig:esagono1}. L'area del triangolo equilatero dato il lato $l$ è \[A_T=\dfrac{l^2}{4}\sqrt{3}\]
	Per quanto detto nel~\vref{thm:areaesagonoin1}
	\[A_E=\dfrac{3}{2}R^2\sqrt{3}\]
	Quindi
	\begin{align*}
		\dfrac{A_E}{A_T}=&\dfrac{\dfrac{3}{2}r^2\sqrt{3}}{\dfrac{l^2}{4}\sqrt{3}}\\
		=&\dfrac{3}{2}r^2\sqrt{3}\dfrac{4}{l^2\sqrt{3}}\\
		=&6\dfrac{r^2}{l^2}\\
		\intertext{Ma per il~\vref{cor:latotriangoloequilatero1}}
		l=&r\sqrt{3}\\
		=&6\dfrac{r^2}{3r^2}\\
		=2
	\end{align*}
	Come volevasi dimostrare.
\end{proof}
\begin{thm}[Apotema dell'esagono]
	in un esagono l'apotema è:\[a=\dfrac{R}{2}\sqrt{3} \]
\end{thm}
\begin{figure}
	\centering
	\includestandalone{geometria/esagono_inscritta_circoscritta}
	\caption{Esagono inscritto e circoscritto}
	\label{fig:esagonoinscrittocircoscritto}
\end{figure}
